\section{Robustness and scope}\label{sec:short_robustness}

Supporting checks address historical specification, numerical resolution, market activity, and external transfer. Online Supplement Section S1 reports all historical benchmarks, Section S2 covers numerical accuracy and posterior calibration, Section S3 assesses liquidity and date dependence, and Section S4 examines LO transfer and tree refinement.

Changing the physical-return model or return source leaves the central comparison intact. Standardized Student-$t$ innovations widen the volatility posterior and modestly improve historical-volatility pricing, while PI and PM remain close. Contract-by-contract reconstruction of the first-nearby CL return history changes the 2,381-row historical MAE/RMSE from 2.6187/3.4090 to 2.6468/3.4451. Thus the historical proxy's roll construction does not explain the observed advantage of option-implied inputs.

Numerical accuracy is important because the PI--PM corrections are small. High-precision Monte Carlo, independently scrambled Sobol pricing, and Curran conditioning recover representative corrections at a scale well above their stochastic numerical noise. They can differ in absolute prices by more than the correction itself, so agreement on PI--PM is checked separately from absolute price agreement. Simulation-based calibration tests the Gaussian posterior implementation. Repricing at posterior RMS volatility, $\sqrt{\mathbb E[\sigma^2\mid\mathcal D]}$, also changes pooled errors only slightly relative to posterior mean volatility under the same Curran approximation.

The LO result is insensitive to the evaluated scalar-transfer alternative. Matching the unresolved average's variance changes previous-day LO RMSE from 0.2599 to 0.2612 on the same 253 rows. An 80/160/320-step CRR audit shows practical stability of the production inversion. These checks support the reported comparison at its numerical scale; they do not establish an exact tree limit or an arbitrage-free fitted surface.

Market activity restricts interpretation. Many targets are end-of-day settlements with zero reported trading volume. The option-implied improvement persists in the pooled positive-volume sample and under open-interest restrictions, but transaction-active observations are sparse within some long-dated expiries. Date-cluster resampling preserves contemporaneous dependence among contracts; it does not eliminate serial dependence across neighboring dates. Fifteen target dates in the main forward exercise and ten in the LO comparison provide limited temporal evidence despite the larger number of contract-date rows.

The pricing model fixes constant volatility and one common futures factor. Implied volatility under this model can absorb risk premia, stochastic volatility, jumps, omitted curve dependence, liquidity, and settlement marking. The comparison therefore concerns the information supplied to a conditional valuation model, rather than an identified structural difference between physical and risk-neutral diffusion coefficients. Richer models could change both absolute fit and the sensitivity of prices to additional uncertain parameters. The present results identify the integration mechanism and its scale within the stated model and samples.
